\subsection{$V_{\mathrm{quad}}$ as the Stokes connection constant of the cubic-coefficient ODE}\label{subsec:vquad-stokes-gcf}

The quadratic-denominator continued fraction whose Borel/PCF value is $\mathcal{V}_{\mathrm{quad}}=1+\mathbf{K}_{n\ge 1}1/(3n^2+n+1)$ is the first genuinely non-Bessel case in the hierarchy discussed above: the same recurrence-to-continuum passage that underlies the super-exponential law~\eqref{eq:wkb} leads to
\[
  (3x^2+x+1)\,y''(x) + (6x+1)\,y'(x) - x^2 y(x) = 0.
\]
Thus the quadratic GCF/Borel problem is accompanied by a genuine Stokes matching problem at the irregular singular point $x=\infty$, and it is natural to isolate the corresponding connection coefficient as an intrinsic analytic invariant of the cubic-coefficient ODE.

To analyze $x=\infty$, insert the WKB ansatz $y(x)\sim x^{\mu}e^{\sigma x}$ into the differential equation. The $x^2$-balance gives $3\sigma^2-1=0$, hence $\sigma_{\pm}=\pm 1/\sqrt{3}$; matching the next $x$-order terms yields $6\sigma\mu+\sigma^2+6\sigma=0$, so
\[
  \mu_{+}=-1-\frac{1}{6\sqrt{3}},
  \qquad
  \mu_{-}=-1+\frac{1}{6\sqrt{3}}.
\]
We therefore define the formal dominant branch (which grows like $e^{x/\sqrt{3}}$ as $x\to+\infty$) and the recessive branch (which decays like $e^{-x/\sqrt{3}}$) on the positive real axis by
\[
  y_{\mathrm{dom}}(x)\sim x^{\mu_+}e^{x/\sqrt{3}}\sum_{n\ge 0} a_n x^{-n},
  \qquad
  y_{\mathrm{rec}}(x)\sim x^{\mu_-}e^{-x/\sqrt{3}}\sum_{n\ge 0} b_n x^{-n},
\]
with $a_0=b_0=1$.

At $x=0$ the indicial roots are $0$ and $1$, so there are no logarithmic terms and one has an analytic Frobenius basis $y_1(x)=1+O(x^4)$, $y_2(x)=x+O(x^2)$ with $y_1(0)=1$. We define the Stokes connection constant $V_{\mathrm{quad}}$ by the coefficient of $y_1$ in the exact local expansion
\[
  y_{\mathrm{rec}}(x)=V_{\mathrm{quad}}\,y_1(x)+C_{\mathrm{quad}}\,y_2(x)
\]
near $x=0$; equivalently, if
\[
  \begin{pmatrix} y_{\mathrm{rec}} & y_{\mathrm{dom}} \\ y'_{\mathrm{rec}} & y'_{\mathrm{dom}} \end{pmatrix}
  =
  \begin{pmatrix} y_1 & y_2 \\ y_1' & y_2' \end{pmatrix} M
\]
at the matching point $x_{\mathrm{match}}=10$, then $V_{\mathrm{quad}}=M_{11}$. Numerically,
\[
  V_{\mathrm{quad}}=M_{11}=1.9420321374711220465\ldots,
\]
while the second Frobenius coefficient satisfies
\[
  M_{21}=-2.9999268666050110215\ldots=-3+\delta,
  \qquad
  \delta\approx 7.313\times 10^{-5},
\]
and the 500-digit rerun confirms that this small deviation from $-3$ is genuine rather than numerical noise. This is distinct from the raw transported sectorial value $y_{\mathrm{rec}}(0)=1.94257302477940314397\ldots$ (with $y'_{\mathrm{rec}}(0)/y_{\mathrm{rec}}(0)=-1.54473580935967544255\ldots$) obtained from the earlier RK4 Stokes integration at dps$=300$: that quantity depends on the normalization of $y_{\mathrm{rec}}$, whereas $V_{\mathrm{quad}}=M_{11}$ is the normalization-independent coefficient of $y_1$.

\begin{theorem}\label{thm:vquad-exclusion}
The Stokes connection constant $V_{\mathrm{quad}}$ does not satisfy any non-trivial integer-linear relation with coefficients bounded by $500$ in the explicit basis
\[
\begin{aligned}
\mathcal{B}=\{&1,\,\pi,\,\log 2,\,\log 3,\,\sqrt{3},\,\Gamma(1/3),\,\Gamma(2/3),\,\Gamma(1/4),\,\Gamma(3/4),\,\zeta(3),\\
&{}_0F_2(;2/3,4/3;4/27),\,{}_0F_2(;1/3,2/3;4/27),\,{}_0F_2(;1/3,4/3;4/27),\\
&{}_0F_2(;2/3,5/3;4/27),\,{}_0F_2(;1/3,2/3;1/27),\,{}_0F_2(;2/3,4/3;1/27),\\
&\operatorname{Ai}\!\left((11/12)^{1/3}\right),\,\operatorname{Bi}\!\left((11/12)^{1/3}\right),\,\operatorname{Ai}\!\left(-(11/12)^{1/3}\right),\,\operatorname{Bi}\!\left(-(11/12)^{1/3}\right),\\
&\operatorname{Ai}(3^{-1/3}),\,\operatorname{Bi}(3^{-1/3}),\,\operatorname{Ai}(11^{1/3}/3),\,\operatorname{Bi}(11^{1/3}/3),\\
&\operatorname{Ai}\!\left((\sqrt{11}/6)^{2/3}\right),\,\operatorname{Bi}\!\left((\sqrt{11}/6)^{2/3}\right),\\
&I_{1/3}\!\left(\tfrac{2}{3\sqrt{3}}\right),\,K_{1/3}\!\left(\tfrac{2}{3\sqrt{3}}\right),\,I_{2/3}\!\left(\tfrac{2}{3\sqrt{3}}\right),\,K_{2/3}\!\left(\tfrac{2}{3\sqrt{3}}\right),\\
&I_{1/3}\!\left(\tfrac{2}{3}\right),\,K_{1/3}\!\left(\tfrac{2}{3}\right)\}.
\end{aligned}
\]
This was verified via the PSLQ algorithm with \texttt{maxcoeff}=500 at \texttt{dps}=300.
\end{theorem}

Structurally, this is the same phenomenon seen for Stokes multipliers in Painlev\'e~II/III and in confluent-Heun problems, where the connection constants are typically more primitive than isolated classical special values. In that sense $V_{\mathrm{quad}}$ behaves less like a disguised Gamma-ratio and more like an entry of a global monodromy/Stokes matrix. The obstruction to a direct ${}_0F_2$ reduction is already visible in the Maclaurin recurrence, whose $c_{n-2}$ term prevents the first-order Pochhammer structure of a single generalized hypergeometric solution. The persistent PSLQ failures therefore point to structure, not to numerical accident. The natural next step is a full symbolic monodromy computation or a connection to the theory of resurgence/alien calculus.
